\chapter{Wasserstein 距離}
\label{ch:seminar-05-wasserstein}

最適輸送の最適値は，地の点の間の距離を「持ち上げて」
測度（ヒストグラム）の間の距離を定める．
本章では，地のコストが距離の $p$ 乗であるとき，
最適輸送が確率測度の空間上の距離
——\textbf{Wasserstein 距離}——を与えることを示す．
証明の核心は\textbf{貼り合わせ補題}（gluing lemma）による三角不等式であり，
これが幾何的発展（測地線・曲率）の出発点となる．

%% ============================================================
%% 5.1 地の距離と Wasserstein 距離
%% ============================================================
\section{地の距離と Wasserstein 距離}
\label{sec:sem-wasserstein-def}

まず離散の設定を固定する．
$\mathbf{a}, \mathbf{b}, \mathbf{c}$ は同一の地の集合 $\range{n}$ 上の
ヒストグラムであり，\textbf{確率単体}
\[
  \simplex_n
  \defeq
  \left\{
    \mathbf{a} \in \R_{\geq 0}^n
    \;\middle|\;
    \textstyle\sum_{i=1}^n a_i = 1
  \right\}
\]
に属するとする．

\begin{definition}{地の距離行列}{sem-ground-metric}
  行列 $\distD \in \R_{\geq 0}^{n \times n}$ が
  $\range{n}$ 上の\textbf{距離（地の距離行列）}であるとは，
  次を満たすことをいう．
  \begin{enumerate}
    \item[\textrm{(i)}] 対称性：$\distD_{i,j} = \distD_{j,i}$．
    \item[\textrm{(ii)}] 非退化性：$\distD_{i,j} = 0 \iff i = j$．
    \item[\textrm{(iii)}] 三角不等式：$\distD_{i,k} \leq \distD_{i,j} + \distD_{j,k}$
      （$\forall\, i, j, k$）．
  \end{enumerate}
  これは距離空間の公理（定義~\ref{def:sem-metric-space}）を
  有限集合 $\range{n}$ に適用したものである．
\end{definition}

\begin{definition}{離散 Wasserstein 距離}{sem-wasserstein-discrete}
  地の距離行列 $\distD$ と $p \geq 1$ に対し，コスト行列を
  $\mathbf{C} = \distD^p = (\distD_{i,j}^p)_{i,j}$ とする．
  $\mathbf{a}, \mathbf{b} \in \simplex_n$ に対して
  \textbf{$p$-Wasserstein 距離}を
  \[
    \WassD_p(\mathbf{a}, \mathbf{b})
    \defeq
    \MKD_{\distD^p}(\mathbf{a}, \mathbf{b})^{1/p}
    =
    \left(
      \min_{\mathbf{P} \in \CouplingsD(\mathbf{a}, \mathbf{b})}
      \inner{\mathbf{P}}{\distD^p}
    \right)^{1/p}
  \]
  で定める（$\WassD_p$ は $\distD$ に依存する）．
\end{definition}

三角不等式の証明には，重み付き $\ell^p$ ノルムの三角不等式
——\textbf{Minkowski の不等式}——を用いる．

\begin{claim}{Minkowski の不等式（重み付き）}{sem-minkowski}
  $p \geq 1$，有限添字集合 $I$，非負の重み $(w_\iota)_{\iota \in I}$，
  実数 $(u_\iota)_{\iota \in I}, (v_\iota)_{\iota \in I}$ に対して，
  \[
    \left(
      \sum_{\iota \in I} w_\iota |u_\iota + v_\iota|^p
    \right)^{1/p}
    \leq
    \left(
      \sum_{\iota \in I} w_\iota |u_\iota|^p
    \right)^{1/p}
    +
    \left(
      \sum_{\iota \in I} w_\iota |v_\iota|^p
    \right)^{1/p}.
  \]
\end{claim}

\begin{proof}
  重み付き内積空間における $L^p$ ノルム
  $\norm{f}_{p,w} \defeq (\sum_\iota w_\iota |f_\iota|^p)^{1/p}$ の
  三角不等式にほかならない．これは古典的な Minkowski の不等式である
  （Hölder の不等式から従う標準的事実なので証明は省く）．
\end{proof}

%% ============================================================
%% 5.2 距離の公理（離散）
%% ============================================================
\section{\texorpdfstring{$\WassD_p$}{Wp} が距離であること}
\label{sec:sem-wasserstein-distance-discrete}

\begin{theorem}{離散 Wasserstein 距離の距離性}{sem-wasserstein-is-distance}
  $\distD$ を $\range{n}$ 上の地の距離行列，$p \geq 1$ とする．
  このとき $\WassD_p$ は確率単体 $\simplex_n$ 上の距離である．
  すなわち，任意の $\mathbf{a}, \mathbf{b}, \mathbf{c} \in \simplex_n$ に対し
  \begin{enumerate}
    \item[\textrm{(i)}] 対称性：$\WassD_p(\mathbf{a}, \mathbf{b})
      = \WassD_p(\mathbf{b}, \mathbf{a})$；
    \item[\textrm{(ii)}] 非退化性：$\WassD_p(\mathbf{a}, \mathbf{b}) \geq 0$ かつ
      $\WassD_p(\mathbf{a}, \mathbf{b}) = 0 \iff \mathbf{a} = \mathbf{b}$；
    \item[\textrm{(iii)}] 三角不等式：$\WassD_p(\mathbf{a}, \mathbf{c})
      \leq \WassD_p(\mathbf{a}, \mathbf{b}) + \WassD_p(\mathbf{b}, \mathbf{c})$．
  \end{enumerate}
\end{theorem}

\begin{proof}
  最適カップリングの存在は主張~\ref{clm:sem-discrete-kantorovich-existence} による．

  \textbf{(i) 対称性．}
  $\distD^p$ は対称なので，転置 $\mathbf{P} \mapsto \mathbf{P}^\top$ は
  $\CouplingsD(\mathbf{a}, \mathbf{b})$ と $\CouplingsD(\mathbf{b}, \mathbf{a})$ の間の
  全単射であり，$\inner{\mathbf{P}^\top}{\distD^p}
  = \inner{\mathbf{P}}{\distD^p}$ で目的関数を保つ．
  ゆえに両者の最小値は等しい．

  \textbf{(ii) 非退化性．}
  $\distD^p \geq \mathbf{0}$ より $\WassD_p \geq 0$．
  $\mathbf{a} = \mathbf{b}$ のとき，$\mathbf{P}^\star = \diag(\mathbf{a})$ は
  $\CouplingsD(\mathbf{a}, \mathbf{a})$ に属し（行和・列和とも $\mathbf{a}$），
  対角が $0$ なので
  $\inner{\diag(\mathbf{a})}{\distD^p} = \sum_i a_i \distD_{i,i}^p = 0$，
  ゆえに $\WassD_p(\mathbf{a}, \mathbf{a}) = 0$．
  逆に $\mathbf{a} \neq \mathbf{b}$ とする．任意の
  $\mathbf{P} \in \CouplingsD(\mathbf{a}, \mathbf{b})$ が対角行列なら
  $\mathbf{a} = \mathbf{P}\ones_n = \mathbf{P}^\top\ones_n = \mathbf{b}$ となり矛盾するので，
  ある $i \neq j$ で $P_{i,j} > 0$．
  非退化性 (ii) より $\distD_{i,j} > 0$ だから
  $\inner{\mathbf{P}}{\distD^p} \geq P_{i,j}\distD_{i,j}^p > 0$．
  最適 $\mathbf{P}$ についてもこれが成り立つので
  $\WassD_p(\mathbf{a}, \mathbf{b}) > 0$．

  \textbf{(iii) 三角不等式（貼り合わせ）．}
  $\mathbf{P} \in \CouplingsD(\mathbf{a}, \mathbf{b})$,
  $\mathbf{Q} \in \CouplingsD(\mathbf{b}, \mathbf{c})$ を，それぞれ
  $(\mathbf{a},\mathbf{b})$, $(\mathbf{b},\mathbf{c})$ の最適カップリングとする．
  $\mathbf{b}$ のゼロ成分を避けるため
  $\tilde{b}_j \defeq b_j$（$b_j > 0$），$\tilde{b}_j \defeq 1$（$b_j = 0$）
  とおき，\textbf{貼り合わせ}行列
  \[
    \mathbf{S} \defeq \mathbf{P}\diag(1/\tilde{\mathbf{b}})\mathbf{Q},
    \qquad
    S_{i,k} = \sum_j \frac{P_{i,j} Q_{j,k}}{\tilde{b}_j}
  \]
  を定める．

  まず $\mathbf{S} \in \CouplingsD(\mathbf{a}, \mathbf{c})$ を示す．
  $b_j = 0$ なら，列 $j$ の和より $\sum_i P_{i,j} = b_j = 0$ かつ
  行 $j$ の和より $\sum_k Q_{j,k} = b_j = 0$ なので，
  その $j$ では $P_{i,j} = Q_{j,k} = 0$．したがって
  \[
    \sum_k S_{i,k}
    = \sum_j \frac{P_{i,j}}{\tilde{b}_j}\sum_k Q_{j,k}
    = \sum_{j : b_j > 0} \frac{P_{i,j}}{b_j} \cdot b_j
    = \sum_j P_{i,j} = a_i,
  \]
  同様に $\sum_i S_{i,k} = c_k$．非負性も明らかなので
  $\mathbf{S} \in \CouplingsD(\mathbf{a}, \mathbf{c})$．

  $\mathbf{S}$ の劣最適性，$\distD$ の三角不等式（および $x \mapsto x^p$ の
  $[0,\infty)$ 上での単調性），重み
  $w_{i,j,k} \defeq P_{i,j}Q_{j,k}/\tilde{b}_j \geq 0$ に対する
  Minkowski の不等式（主張~\ref{clm:sem-minkowski}）を順に用いると，
  \begin{align*}
    \WassD_p(\mathbf{a}, \mathbf{c})
    &\leq \inner{\mathbf{S}}{\distD^p}^{1/p}
     = \left(\sum_{i,j,k} \distD_{i,k}^p\, w_{i,j,k}\right)^{1/p} \\
    &\leq \left(\sum_{i,j,k}
        (\distD_{i,j} + \distD_{j,k})^p\, w_{i,j,k}\right)^{1/p} \\
    &\leq \left(\sum_{i,j,k} \distD_{i,j}^p\, w_{i,j,k}\right)^{1/p}
       + \left(\sum_{i,j,k} \distD_{j,k}^p\, w_{i,j,k}\right)^{1/p}.
  \end{align*}
  第 1 項では $\sum_k Q_{j,k} = b_j$ より
  \[
    \sum_{i,j,k} \distD_{i,j}^p\, \frac{P_{i,j}Q_{j,k}}{\tilde{b}_j}
    = \sum_{i,j} \distD_{i,j}^p\, \frac{P_{i,j}}{\tilde{b}_j} b_j
    = \sum_{i,j} \distD_{i,j}^p\, P_{i,j}
    = \inner{\mathbf{P}}{\distD^p}
    = \WassD_p(\mathbf{a}, \mathbf{b})^p,
  \]
  第 2 項では $\sum_i P_{i,j} = b_j$ より同様に
  $\inner{\mathbf{Q}}{\distD^p} = \WassD_p(\mathbf{b}, \mathbf{c})^p$ を得る．
  （いずれも $b_j = 0$ の項は $P_{i,j} = Q_{j,k} = 0$ により寄与しない．）
  よって
  $\WassD_p(\mathbf{a}, \mathbf{c})
  \leq \WassD_p(\mathbf{a}, \mathbf{b}) + \WassD_p(\mathbf{b}, \mathbf{c})$．
\end{proof}

\begin{remark}{$0 < p \leq 1$ の場合}{sem-wasserstein-p-leq-1}
  $0 < p \leq 1$ のときは $\distD^p$ 自身が地の距離行列となり，
  $\WassD_p(\mathbf{a}, \mathbf{b})^p$（$1/p$ 乗を外したもの）が
  $\simplex_n$ 上の距離となる．以下では $p \geq 1$ を仮定する．
\end{remark}

%% ============================================================
%% 5.3 連続版への一般化
%% ============================================================
\section{連続版への一般化}
\label{sec:sem-wasserstein-continuous}

同じ構成は任意の確率測度に拡張される．
$\X$ を Polish 空間，$d$ を $\X$ 上の距離とし，
$p$ 次モーメントが有限な確率測度の空間
\[
  \Pp_p(\X)
  \defeq
  \left\{
    \alpha \in \Mm_+^1(\X)
    \;\middle|\;
    \int_\X d(x, x_0)^p \,\d\alpha(x) < \infty
    \ (\exists\, x_0 \in \X)
  \right\}
\]
を考える（有限性は基準点 $x_0$ の取り方によらない）．

\begin{definition}{Wasserstein 距離}{sem-wasserstein-cont-def}
  $p \geq 1$ と $\alpha, \beta \in \Pp_p(\X)$ に対して，
  \textbf{$p$-Wasserstein 距離}を
  \[
    \Wass_p(\alpha, \beta)
    \defeq
    \MK_{d^p}(\alpha, \beta)^{1/p}
    =
    \left(
      \inf_{\pi \in \Couplings(\alpha, \beta)}
      \int_{\X \times \X} d(x, y)^p \,\d\pi(x, y)
    \right)^{1/p}
  \]
  で定める．
\end{definition}

連続版の三角不等式には，貼り合わせ行列
$\mathbf{P}\diag(1/\tilde{\mathbf{b}})\mathbf{Q}$ の連続版が必要となる．

\begin{claim}{貼り合わせ補題}{sem-gluing-lemma}
  $\alpha, \beta, \gamma \in \Mm_+^1(\X)$ とし，
  $\pi \in \Couplings(\alpha, \beta)$,
  $\sigma \in \Couplings(\beta, \gamma)$ とする．
  このとき $\X^3$ 上の確率測度 $\mu$ が存在して，
  $(x,y)$ 成分の周辺が $\pi$，$(y,z)$ 成分の周辺が $\sigma$ となる．
  とくに $(x,z)$ 成分の周辺
  $\mathbf{S} \defeq (P_{x,z})_\sharp \mu$ は
  $\Couplings(\alpha, \gamma)$ に属する．
\end{claim}

\begin{proof}[証明の方針]
  $\beta$ に関して $\pi, \sigma$ を分解（disintegration）し，
  $y$ で条件づけた測度 $\pi_y, \sigma_y$ を用いて
  $\d\mu(x,y,z) = \d\beta(y)\,\d\pi_y(x)\,\d\sigma_y(z)$ とおけばよい．
  これは離散の $S_{i,k} = \sum_j P_{i,j}Q_{j,k}/\tilde{b}_j$ の連続版である．
  分解定理（Polish 空間上で成立）の詳細は測度論の標準的結果に譲る．
\end{proof}

\begin{proposition}{Wasserstein 距離の距離性}{sem-wasserstein-cont-distance}
  $d$ が $\X$ 上の距離，$p \geq 1$ のとき，
  $\Wass_p$ は $\Pp_p(\X)$ 上の距離である．
\end{proposition}

\begin{proof}
  対称性と非退化性は離散の場合（定理~\ref{thm:sem-wasserstein-is-distance}）と
  同様に，$d^p$ の対称性と，対角線
  $\{(x,x)\}$ に台を持つカップリング（$\alpha = \beta$ のとき）の利用，
  および $d(x,y) > 0$（$x \neq y$）から従う．

  三角不等式を示す．$\pi \in \Couplings(\alpha, \beta)$,
  $\sigma \in \Couplings(\beta, \gamma)$ を最適カップリングとし，
  主張~\ref{clm:sem-gluing-lemma} の $\mu$ をとる．
  $\mathbf{S} = (P_{x,z})_\sharp \mu \in \Couplings(\alpha, \gamma)$ の
  劣最適性，$d$ の三角不等式，測度 $\mu$ に関する Minkowski の不等式より
  \begin{align*}
    \Wass_p(\alpha, \gamma)
    &\leq \left(\int d(x,z)^p \,\d\mu\right)^{1/p}
     \leq \left(\int \bigl(d(x,y) + d(y,z)\bigr)^p \,\d\mu\right)^{1/p} \\
    &\leq \left(\int d(x,y)^p \,\d\mu\right)^{1/p}
       + \left(\int d(y,z)^p \,\d\mu\right)^{1/p} \\
    &= \left(\int d(x,y)^p \,\d\pi\right)^{1/p}
       + \left(\int d(y,z)^p \,\d\sigma\right)^{1/p}
     = \Wass_p(\alpha, \beta) + \Wass_p(\beta, \gamma).
  \end{align*}
  最後の等号は $\mu$ の $(x,y)$ 周辺が $\pi$，$(y,z)$ 周辺が $\sigma$ で
  あることによる．
\end{proof}

\begin{remark}{離散版との整合}{sem-wasserstein-disc-cont}
  ヒストグラム $\mathbf{a} = \sum_i a_i\delta_{x_i}$,
  $\mathbf{b} = \sum_j b_j\delta_{x_j}$ を同一の台
  $\{x_1, \ldots, x_n\} \subset \X$ 上の離散測度と見なすと，
  $\distD_{i,j} = d(x_i, x_j)$ のもとで
  $\Wass_p$ は $\WassD_p$ に一致する．
  定理~\ref{thm:sem-wasserstein-is-distance} は
  命題~\ref{prop:sem-wasserstein-cont-distance} の離散特別な場合である．
\end{remark}

%% ============================================================
%% 5.4 弱距離としての Wasserstein 距離
%% ============================================================
\section{弱距離としての Wasserstein 距離}
\label{sec:sem-wasserstein-weak}

Wasserstein 距離の重要な特長は，台が重ならない測度
（たとえば離散測度）の間の「ずれ」を測れる\textbf{弱距離}である点である．
これは古典的な $L^2$ ノルムなどにはない性質である．

\begin{proposition}{Dirac 測度間の Wasserstein 距離}{sem-wasserstein-dirac}
  任意の $x, y \in \X$ と $p \geq 1$ に対して
  \[
    \Wass_p(\delta_x, \delta_y) = d(x, y).
  \]
\end{proposition}

\begin{proof}
  $\Couplings(\delta_x, \delta_y)$ の唯一の元は積測度
  $\delta_x \otimes \delta_y = \delta_{(x,y)}$ である
  （周辺が $\delta_x, \delta_y$ となる測度は $(x,y)$ に全質量を置くものに限る）．
  ゆえに
  $\Wass_p(\delta_x, \delta_y)^p
  = \int d^p \,\d\delta_{(x,y)} = d(x,y)^p$．
\end{proof}

とくに $x \to y$ のとき $\Wass_p(\delta_x, \delta_y) \to 0$ となり，
Wasserstein 距離は台の位置の連続的な移動を捉える．
これを一般化するのが\textbf{弱収束}である．

\begin{definition}{測度の弱収束}{sem-weak-convergence}
  $\X$ 上の有界連続関数全体を $\Cc_b(\X)$ と書く．
  $\Mm_+^1(\X)$ の列 $(\alpha_k)_k$ が $\alpha$ に\textbf{弱収束}するとは，
  任意の $g \in \Cc_b(\X)$ に対して
  \[
    \int_\X g \,\d\alpha_k \;\longrightarrow\; \int_\X g \,\d\alpha
    \qquad (k \to \infty)
  \]
  が成り立つことをいい，$\alpha_k \rightharpoonup \alpha$ と書く．
  これは確率変数の分布収束に対応する．
\end{definition}

\begin{claim}{Wasserstein 距離は弱収束を距離化する}{sem-wasserstein-metrizes-weak}
  $\alpha_k, \alpha \in \Pp_p(\X)$ とする．このとき
  \[
    \Wass_p(\alpha_k, \alpha) \to 0
    \iff
    \alpha_k \rightharpoonup \alpha
    \ \text{かつ}\
    \int d(x, x_0)^p \,\d\alpha_k \to \int d(x, x_0)^p \,\d\alpha.
  \]
  すなわち $\Wass_p$ は，$p$ 次モーメントの収束を伴う弱収束を距離化する．
  とくに $\X$ が有界距離空間ならモーメント条件は自動的に成り立ち，
  $\Wass_p$ 収束は弱収束と同値である．
\end{claim}

\begin{proof}[証明の方針]
  有界距離空間では右辺のモーメント条件は不要となり，
  $\Wass_p$ 収束と弱収束が同値になる．
  一般の場合の証明は測度の弱位相と
  輸送コストの下半連続性を用いる（標準的結果に譲る）．
\end{proof}

\begin{remark}{本章のまとめと展望}{sem-ch5-summary}
  地のコストが距離の $p$ 乗のとき，最適輸送は確率測度の空間
  $\Pp_p(\X)$ 上の距離 $\Wass_p$ を定める
  （定理~\ref{thm:sem-wasserstein-is-distance}，
  命題~\ref{prop:sem-wasserstein-cont-distance}）．
  これにより $(\Pp_p(\X), \Wass_p)$ は距離空間となり，
  その中で「測地線」や「曲率」を論じる舞台が整う．
  測地線を具体的に構成するには，最適カップリングが写像
  $T = \nabla\phi$ に台を持つという Brenier の定理が必要であり，
  そのためにまず非正則化 Kantorovich 問題の双対と
  $c$-変換を整備する．これが次章以降の主題である．
\end{remark}
